% !PRINTED 175
\weberpage{175}
If in these equations one expresses $\omega'$ in terms of $\omega$ according to (9), one obtains

\[
\begin{aligned}
1 &= \alpha'(\alpha+\beta\omega) + \beta'(\gamma+\delta\omega),\\
\omega &= \gamma'(\alpha+\beta\omega) + \delta'(\gamma+\delta\omega),
\end{aligned}
\]

and since $\omega$ is not real, each of these equations splits into two others:

\[
\begin{aligned}
\alpha\alpha' + \gamma\beta' &= 1, & \beta\alpha' + \delta\beta' &= 0,\\
\alpha\gamma' + \gamma\delta' &= 0, & \beta\gamma' + \delta\delta' &= 1.
\end{aligned}
\]

Consequently

\[
(\alpha\delta - \beta\gamma)(\alpha'\delta' - \beta'\gamma') = 1,
\]

and therefore, since $\omega,\omega'$ must both have a positive imaginary part, so that by (9) the determinant $\alpha\delta-\beta\gamma$ must be positive:

\[
\alpha\delta - \beta\gamma = 1,
\]
\[
\alpha = \delta', \quad \delta = \alpha', \quad \gamma = -\gamma', \quad \beta = -\beta'.
\]

1. \gesperrt{Thus $\omega, \omega'$ are connected with one another by a linear transformation, and indeed [by (6), (8)] by one of the first class (\S\ 36). That conversely two such values $\omega, \omega'$ also give the same value $\varkappa^2$ has already been shown above.}

A similar theorem follows from this as an immediate consequence, for the invariant $j(\omega)$ (\S\ 46).

2. \gesperrt{The equation:}

\weq{10}{j(\omega') = j(\omega)}

\gesperrt{is then and} \textbf{only} \gesperrt{then satisfied, when $\omega'$ is connected with $\omega$ by some linear transformation}

\[
\begin{pmatrix} \alpha, & \beta \\ \gamma, & \delta \end{pmatrix}
\]

\gesperrt{that is, when}

\weq{11}{\omega' = \frac{\gamma+\delta\omega}{\alpha+\beta\omega}, \quad \alpha\delta-\beta\gamma = 1}

\gesperrt{holds.}

For if we denote the values of $\varkappa^2,\varkappa'^2$ belonging to $\omega'$ by $\lambda^2,\lambda'^2$, then equation (10) can be written thus:

\[
\frac{(1-\lambda^2\lambda'^2)^3}{\lambda^4\lambda'^4} = \frac{(1-\varkappa^2\varkappa'^2)^3}{\varkappa^4\varkappa'^4},
\]

and with respect to $\lambda^2$ it is an equation of the sixth degree, whose roots

\[
\varkappa^2, \quad \varkappa'^2, \quad \frac{1}{\varkappa^2}, \quad \frac{1}{\varkappa'^2}, \quad -\frac{\varkappa'^2}{\varkappa^2}, \quad -\frac{\varkappa^2}{\varkappa'^2}
\]
%CONT>>
